\documentclass[conference]{IEEEtran}
\IEEEoverridecommandlockouts

% --- Typography & Core Packages ---
\usepackage{fontspec}
\setmainfont{TeX Gyre Termes}
\usepackage{cite}
\usepackage{amsmath,amssymb,amsfonts}
\usepackage{algorithmic}
\usepackage{graphicx}
\usepackage{textcomp}
\usepackage{xcolor}
\usepackage{booktabs}
\usepackage{tabularx}
\usepackage{url}
\usepackage{hyperref}
\hypersetup{colorlinks=true, linkcolor=blue, citecolor=blue, urlcolor=blue}
\usepackage{multirow}
\usepackage{makecell}

% --- Auto-wrapped column ---
\newcolumntype{Y}{>{\raggedright\arraybackslash}X}

\begin{document}

\title{Reading Guide to the DROS Trilogy: An Agent Runtime Operation Substrate\\
{\footnotesize \textsuperscript{*}Technical Note: Architectural Taxonomy, Enforcement Boundaries, and Cross-Paper Navigation}}

\author{
\IEEEauthorblockN{Chun-Cheng (Jimmy) Chen}
\IEEEauthorblockA{\textit{Top-Celestial Company Ltd.} \\
Taipei, Taiwan R.O.C. \\
jimmychen@dr-os.io}
}

\maketitle
\thispagestyle{plain}
\pagestyle{plain}

\begin{abstract}
As autonomous AI agents are increasingly entrusted with effectful execution in enterprise workflows and critical infrastructure, conventional defenses reliant solely on probabilistic semantic guardrails or coarse-grained operating system sandboxes fail to address the fundamental \textit{Agent-to-Execution Attribution Gap}. The Dharma/Deterministic Runtime Operation Substrate (DROS) research series resolves this architectural deficiency. This Technical Note serves as an authoritative reading guide and taxonomy navigator across the DROS Academic Trilogy: (1) \textbf{DROS-6P} establishes the minimal closed-loop semantic specification across six essential enterprise trust boundaries; (2) \textbf{DROS 4-Layer (v3)} formalizes the attribution gap and implements an in-band sub-microsecond defense-in-depth funnel; and (3) \textbf{DROS-PGM} delivers an unbypassable, kernel-anchored deterministic execution control plane under post-compromise threat models. We explicitly define the exact enforcement scopes, clarify the epistemic boundaries of empirical benchmarks, and provide a unified citation topology for peer reviewers, security researchers, and systems architects.
\end{abstract}

\begin{IEEEkeywords}
AI Agent Governance, Runtime Operation Substrate, Execution Boundary Enforcement, Post-Compromise Security, Defense-in-Depth, Capability Bitmap.
\end{IEEEkeywords}

\section{Purpose and Architectural Role}
The DROS trilogy does not represent disparate point solutions; rather, it constitutes a unified, systems-level convergence addressing three distinct operational facets:
\begin{itemize}
    \item \textbf{What to Govern (Completeness):} Defining the minimal closed-loop trust services necessary for an autonomous agent to act as a legally accountable entity.
    \item \textbf{Where and How to Enforce (Executability):} Translating high-level governance policies into an in-band, lock-free, sub-microsecond runtime control plane.
    \item \textbf{Containment Post-Compromise (Resilience):} Guaranteeing unbypassable execution containment even when user-space agent processes or credentials are fully subverted.
\end{itemize}

\begin{quote}
\textbf{Core Definition:} \textit{DROS is a deterministic runtime governance substrate for autonomous AI workloads, providing an enforceable control boundary between agent-originated intent and effectful execution.}
\end{quote}

By analogy to POSIX in general-purpose computing, POSIX does not eliminate application-level business complexity, but standardizes the low-level execution interface so that every program need not reimplement system calls. DROS adopts this exact abstraction principle: upper-level agent frameworks and multi-agent meshes need not reinvent identity attestation, capability bitmaps, tool mediation, policy gates, cryptographic audit trails, and instantaneous atomic revocation.

\section{Trilogy Overview and Core Papers}
The structural division, core research questions, and authoritative Zenodo DOI citations of the trilogy are summarized in Table~\ref{tab:trilogy_overview_en}.

\begin{table*}[t]
\caption{DROS Trilogy Architecture, Core Research Questions, and Authoritative DOI Registry}
\label{tab:trilogy_overview_en}
\centering
\footnotesize
\renewcommand{\arraystretch}{1.3}
\begin{tabularx}{\textwidth}{l Y Y Y}
\toprule
\textbf{Paper Identifier} & \textbf{Core Research Question} & \textbf{System Architectural Role} & \textbf{Authoritative DOI Citation} \\
\midrule
\textbf{Paper 1: DROS-6P} & Can enterprise AI agent compliance be deterministically enforced across a closed-loop runtime service specification? & \textbf{Requirements \& Semantic Specification}: Principal, Authorization, Tool Bound, Policy Gate, Audit Log, Expiry/Revocation (6P). & DOI: \href{https://doi.org/10.5281/zenodo.21833970}{10.5281/zenodo.21833970} \\
\textbf{Paper 2: DROS 4-Layer (v3)} & Why are semantic firewalls or OS sandboxes insufficient alone? How to bridge the Agent-to-Execution Attribution Gap? & \textbf{Defense-in-Depth Funnel}: L1$\rightarrow$L4 layers; L4 constant-time capability bitmap enforcement at C-ABI/FFI boundaries; ablation \& AGT comparative study. & DOI: \href{https://doi.org/10.5281/zenodo.22092008}{10.5281/zenodo.22092008} \\
\textbf{Paper 3: DROS-PGM} & When credentials expire or user-space processes are fully compromised, can the substrate maintain deterministic containment? & \textbf{Post-Compromise Execution Trust}: Three-plane decoupling, sub-microsecond in-band physical fusing, and formal proof of Unbypassable Mediation. & DOI: \href{https://doi.org/10.5281/zenodo.21903687}{10.5281/zenodo.21903687} \\
\bottomrule
\end{tabularx}
\end{table*}

\textbf{Recommended Reading Path:}
\begin{center}
\textbf{Paper 1 (DROS-6P)} $\longrightarrow$ \textbf{Paper 2 (DROS 4-Layer v3)} $\longrightarrow$ \textbf{Paper 3 (DROS-PGM)}
\end{center}
Readers are encouraged to first establish ``what must be governed'', proceed to ``how enforcement occurs at the execution boundary'', and conclude with ``how containment holds post-compromise''.

\section{Macro Topology: They Decide, DROS Enforces}
A key theoretical contribution of the trilogy is clarifying the boundary between application-layer governance middleware (e.g., Microsoft AGT, LangChain, MCP Policy SDKs) and execution-boundary enforcement (Fig.~\ref{fig:topology_en}).

\begin{figure}[!t]
\centering
\begin{verbatim}
       +------------------------------------+
       |  AI Agent Applications / Workflows |
       +------------------------------------+
              |                      |
           Ingress                Egress
       (Data/Prompt Entry)    (Effectful Action/Tool)
              |                      |
              v                      v
       +------------------------------------+
       |   DROS Runtime Operation Substrate |
       |   - 6P Closed-Loop Service Specs   |
       |   - L1-L4 Probabilistic-to-        |
       |     Deterministic Defense Funnel   |
       |   - Post-Compromise PGM Engine     |
       +------------------------------------+
                         |
                 Deterministic ALLOW
                         |
                         v
       +------------------------------------+
       |  OS Kernel / Storage / Network APIs|
       +------------------------------------+
\end{verbatim}
\caption{DROS Runtime Governance Substrate Topology and Execution Boundaries}
\label{fig:topology_en}
\end{figure}

Paper 2 (v3) succinctly formulates this complementary relationship:
\begin{quote}
\textit{``They decide. DROS enforces.''} \\
(Application middleware excels at declarative policy reasoning and workflow orchestration; unmanaged native execution paths require deterministic C-ABI/FFI physical boundary enforcement.)
\end{quote}

\section{The 6P Closed-Loop Service Matrix}
The six pillars established in DROS-6P constitute the minimal operational baseline for autonomous agents operating in regulated environments (Table~\ref{tab:6p_matrix_en}).

\begin{table}[h]
\caption{6P Runtime Service Specification Matrix}
\label{tab:6p_matrix_en}
\centering
\footnotesize
\renewcommand{\arraystretch}{1.3}
\begin{tabularx}{\columnwidth}{l Y Y}
\toprule
\textbf{Dimension} & \textbf{Core Question} & \textbf{Underlying Mechanism} \\
\midrule
\textbf{Principal} & On whose behalf is it acting? & Cryptographic Execution Token (DIT) \& W3C DID \\
\textbf{Authorization} & What actions are permitted? & Compact Capability Bitmap Vector ($O(1)$) \\
\textbf{Tool Bound} & Which invocations can exit? & In-band C-ABI / FFI Sub-microsecond Panic \\
\textbf{Policy Gate} & How are high risks mitigated? & Dynamic Redaction, HITL Multi-signature \\
\textbf{Audit Log} & How to ensure non-repudiation? & SHA-256 Merkle Hash Chain \\
\textbf{Revocation} & How to invalidate immediately? & Lock-Free RCU Atomic Pointer Hot-Swap \\
\bottomrule
\end{tabularx}
\end{table}

\section{Empirical Benchmarks \& Reproducibility}
All empirical metrics reported across the trilogy are tethered to reproducible test suites in the open-source \href{https://github.com/Top-Celestial-Company-Ltd/DROS-VEP-lite}{DROS-VEP-lite} harness:
\begin{enumerate}
    \item \textbf{Latency Decomposition:} High-level policy evaluation latency (e.g., $26.21\,\mu\mathrm{s}$ P50) and low-level C-ABI physical circuit-breaker latency ($<500\,\mathrm{ns}$) measure distinct code paths and must not be conflated.
    \item \textbf{24-Hour Continuous Soak Test:} Across 160,611 adversarial requests, DROS exhibited zero memory leaks (0 Bytes) and $<1.8\%$ CPU overhead, validating the zero-heap-allocation C-ABI microkernel design.
    \item \textbf{Epistemic Claim Calibration:} Statements such as ``100\% adversarial interception'' reflect empirical results under the evaluated threat scenarios and corpus, and do not constitute an unconstrained claim against arbitrary kernel compromises.
\end{enumerate}

\section{Crucial Defensive Scope \& Boundaries}
To maintain rigorous scientific fidelity, we explicitly demarcate the guarantees provided by DROS:
\begin{itemize}
    \item \textbf{Authorization Confinement vs. Semantic Correctness:} DROS guarantees that unauthorized system effects cannot occur; it does not claim that generative models will never output logically flawed or factually incorrect semantic tokens. The AI may err, but its errors cannot manifest as unauthorized system state transitions.
    \item \textbf{TCB Boundaries:} The Trusted Computing Base assumes integrity of the OS kernel, hardware TPM/HSM, and PGM kernel drivers. Compromises entirely subverting the underlying hypervisor or host hardware reside outside this boundary.
    \item \textbf{Complement to Organizational Governance:} Cryptographic Merkle logs provide non-repudiable forensic proof, but must integrate with organizational compliance controls.
\end{itemize}

\section{Conclusion}
The DROS research series converges agent runtime governance into an intelligible, verifiable, and deployable substrate: \textbf{6P defines service completeness, the 4-Layer funnel defines deterministic execution enforcement, and PGM extends identical execution trust into post-compromise environments.}

\begin{quote}
\textit{``Enabling autonomous AI agents to operate safely requires not just more compliant models, but an unbypassable deterministic governance substrate across effectful execution boundaries.''}
\end{quote}
That substrate is what the DROS Trilogy defines and validates.

\section*{Acknowledgment \& AI Collaboration Disclosure}
In accordance with IEEE 2024+ authorship guidelines and enterprise AI governance transparency principles, the author declares that the novel security concepts, physical-layer architecture, 6-pillar framework, mathematical formalisms, and patent claims presented in this work were independently conceptualized, architected, and validated by Chun-Cheng (Jimmy) Chen. Large language models were utilized strictly for textual translation, grammar refinement, and LaTeX typesetting assistance.

\section*{Patent Protection Notice}
DROS execution governance and security technology is protected under U.S. Provisional Patent Application (U.S. PPA No. 64/111,973, Patent Pending).

\bibliographystyle{IEEEtran}
\begin{thebibliography}{00}
\bibitem{dros6p} C.-C. Chen, ``DROS-6P: A Unified Deterministic Runtime Governance Architecture Closing the Six Fundamental Trust Boundaries of Enterprise AI Agents,'' \textit{Zenodo Technical Report}, 2026. DOI: \href{https://doi.org/10.5281/zenodo.21833970}{10.5281/zenodo.21833970}.
\bibitem{dros4l} C.-C. Chen, ``Bridging the Agent-to-Execution Attribution Gap in Autonomous AI Workloads: A 4-Layer Deterministic Runtime Operating System,'' \textit{ACM Transactions on Privacy and Security (Submitted / Zenodo Preprint v3)}, 2026. DOI: \href{https://doi.org/10.5281/zenodo.22092008}{10.5281/zenodo.22092008}.
\bibitem{drospgm} C.-C. Chen, ``DROS-PGM: A Deterministic Kernel-Level Execution Control Plane for Post-Compromise Security in Autonomous AI Systems,'' \textit{Zenodo Technical Report}, 2026. DOI: \href{https://doi.org/10.5281/zenodo.21903687}{10.5281/zenodo.21903687}.
\bibitem{veplite} Top-Celestial Company Ltd., ``DROS Virtual Enterprise Platform Lite (DROS-VEP-lite),'' GitHub Repository, 2026. \url{https://github.com/Top-Celestial-Company-Ltd/DROS-VEP-lite}.
\end{thebibliography}

\end{document}
